\documentclass[11pt]{article}
\usepackage[T1]{fontenc}
\usepackage{textcomp}
\usepackage[margin=0.75in]{geometry}
\usepackage{amsmath,amssymb,amsthm}
\usepackage{array,booktabs}
\usepackage{hyperref}
\setlength{\parindent}{0pt}
\newtheorem{theorem}{Theorem}
\theoremstyle{definition}
\newtheorem{definition}{Definition}
\title{Flow Proof Artifact: sine-derivatives-at-zero}
\author{Generated by \texttt{flow doc proof}}
\date{Flow Proof Book}
\begin{document}
\maketitle
The derivatives of sine at zero follow the alternating Maclaurin pattern 0, 1, 0, -1, …\\[0.5em]
\textbf{Source.} brook taylor — https://en.wikipedia.org/wiki/Taylor\_series\\[0.5em]
\bigskip
\noindent{\large \textbf{Definition 1.} \textit{The derivatives of sine at zero follow the alternating Maclaurin pattern 0, 1, 0, -1, …} \hfill real analysis \cdot smooth functions \cdot derivatives of sine are known}\par
\smallskip\noindent\textit{Coordinate: real analysis \cdot smooth functions \cdot derivatives of sine are known}\par
\smallskip\noindent\textit{Source: brook taylor — https://en.wikipedia.org/wiki/Taylor\_series}\par
\medskip
\noindent\textbf{Goal.} The derivatives of sine at zero follow the alternating Maclaurin pattern 0, 1, 0, -1, …\par
\noindent$\displaystyle \forall x \in \mathbb{R}\quad \sin''(0) = 0$\par
\medskip
\noindent
\begin{tabular}{@{} >{\raggedright\arraybackslash}p{0.06\textwidth} >{\raggedright\arraybackslash}p{0.47\textwidth} >{\raggedleft\arraybackslash}p{0.41\textwidth} @{}}
\textbf{\#} & \textbf{Proof} & \textbf{Mathematics} \\
\midrule
\textbf{1.} & We stipulate derivatives of sine are known for smooth functions on real analysis — this is a definition, not a derived fact. & \\[0.45em]
\textbf{2.} & This follows directly from the definition: the first derivative of sine at zero equals 1. & $\displaystyle \sin'(0) = 1$ \\[0.45em]
\textbf{3.} & This follows directly from the definition: the second derivative of sine at zero equals 0. Hence proven. & $\displaystyle \sin''(0) = 0$ \\[0.45em]
\bottomrule
\end{tabular}
\label{thm:Analysis-smooth functions-derivatives_of_sine_are_known}
\par\medskip\hrule
\end{document}
